% ============================================================
\chapter{Arithm\'etique Flottante, Erreurs et Stabilit\'e}
\label{ch:flottants}
% ============================================================

\begin{center}
\textit{<<~Le calcul num\'erique est l'art de donner la bonne r\'eponse \`a un probl\`eme l\'eg\`erement diff\'erent.~>>}
\end{center}

\section{Exemple motivant}

Calculons $f(x)=\frac{1-\cos x}{x^2}$ pour $x=10^{-8}$ en double pr\'ecision. Le r\'esultat th\'eorique est $\approx 0{,}5$, mais Python donne $0{,}0$ ! La \textbf{cancellation catastrophique} entre $1$ et $\cos x$ \'elimine tous les chiffres significatifs. Comprendre l'arithm\'etique flottante est indispensable.

\section{Repr\'esentation des nombres en machine}

\begin{definition}[Nombre flottant]\index{nombre flottant}
En base $\beta$ avec $t$ chiffres significatifs, un nombre flottant s'\'ecrit :
\[
x = \pm m \times \beta^e, \quad m = d_0.d_1d_2\cdots d_{t-1}, \quad d_0\neq 0,
\]
o\`u $e_{\min}\leq e\leq e_{\max}$ et $0\leq d_i<\beta$.
\end{definition}

\begin{definition}[Standard IEEE 754]\index{IEEE 754}
\begin{center}
\begin{tabular}{lccc}
\toprule
\textbf{Format} & \textbf{Bits} & \textbf{Mantisse} & \textbf{$\eps_{\mathrm{mach}}$} \\
\midrule
Simple pr\'ecision & 32 & 23+1 bits & $\approx 1{,}19\times 10^{-7}$ \\
Double pr\'ecision & 64 & 52+1 bits & $\approx 2{,}22\times 10^{-16}$ \\
\bottomrule
\end{tabular}
\end{center}
\end{definition}

\begin{definition}[Epsilon machine]\index{epsilon machine}
L'\textbf{epsilon machine} $\eps_{\mathrm{mach}}$ est le plus petit $\eps>0$ tel que $\fl(1+\eps)>1$ :
\[
\eps_{\mathrm{mach}} = \frac{1}{2}\beta^{1-t}.
\]
Pour tout $x\in\R$ repr\'esentable, $\fl(x)=x(1+\delta)$ avec $|\delta|\leq\eps_{\mathrm{mach}}$.
\end{definition}

\section{Erreurs d'arrondi et propagation}

\begin{definition}[Erreur absolue et relative]\index{erreur absolue}\index{erreur relative}
\begin{align}
\text{Erreur absolue} &: E_a = |x-\tilde{x}|, \\
\text{Erreur relative} &: E_r = \frac{|x-\tilde{x}|}{|x|} \quad (x\neq 0).
\end{align}
\end{definition}

\begin{theoreme}[Propagation des erreurs]
Si $\tilde{x}=x(1+\eps_x)$ et $\tilde{y}=y(1+\eps_y)$ :
\begin{itemize}
\item $\tilde{x}\cdot\tilde{y}\approx xy(1+\eps_x+\eps_y)$ : erreurs relatives s'ajoutent.
\item $\tilde{x}-\tilde{y}\approx(x-y)+x\eps_x-y\eps_y$ : si $x\approx y$, l'erreur relative explose.
\end{itemize}
\end{theoreme}

\begin{attention}
La \textbf{cancellation catastrophique} se produit lors de la soustraction de nombres presque \'egaux. Exemple : $1-\cos x$ pour $x$ petit. Rem\`ede : utiliser $2\sin^2(x/2)$.
\end{attention}

\section{Conditionnement d'un probl\`eme}

\begin{definition}[Conditionnement]\index{conditionnement}
Le \textbf{conditionnement} d'un probl\`eme $f$ au point $x$ est :
\[
\kappa(f,x) = \left|\frac{xf'(x)}{f(x)}\right|.
\]
Un probl\`eme est \textbf{bien conditionn\'e} si $\kappa\sim 1$, \textbf{mal conditionn\'e} si $\kappa\gg 1$.
\end{definition}

\begin{definition}[Conditionnement matriciel]\index{conditionnement matriciel}
Pour $Ax=b$ :
\[
\cond(A) = \|A\|\cdot\|A^{-1}\|.
\]
L'erreur relative sur $x$ est amplifi\'ee par $\cond(A)$ :
\[
\frac{\|\delta x\|}{\|x\|} \leq \cond(A)\frac{\|\delta b\|}{\|b\|}.
\]
\end{definition}

\section{Stabilit\'e num\'erique}

\begin{definition}[Algorithme stable]\index{stabilite@stabilit\'e}
Un algorithme est \textbf{stable} s'il donne la solution exacte d'un probl\`eme \textbf{l\'eg\`erement perturb\'e}. Il est \textbf{r\'etrograde stable} (backward stable) si la perturbation est de l'ordre de $\eps_{\mathrm{mach}}$.
\end{definition}

\begin{methode}
La pr\'ecision du r\'esultat d\'epend de deux facteurs :
\begin{enumerate}
\item Le \textbf{conditionnement} du probl\`eme (inh\'erent, on ne peut le changer).
\item La \textbf{stabilit\'e} de l'algorithme (choix du programmeur).
\end{enumerate}
R\`egle : erreur $\lesssim \cond \times \eps_{\mathrm{mach}}$.
\end{methode}

\begin{center}
\begin{tikzpicture}[
  box/.style={rectangle, draw, rounded corners, minimum width=2.5cm, minimum height=1.2cm, text centered, font=\small},
  arr/.style={-{Stealth}, thick}
]
\node[box,fill=blue!10,align=center] (prob) at (0,0) {Probl\`eme\\$f(x)$};
\node[box,fill=green!10,align=center] (algo) at (5,0) {Algorithme\\$\tilde{f}(\tilde{x})$};
\node[box,fill=orange!10,align=center] (res) at (10,0) {R\'esultat\\$\tilde{y}$};
\draw[arr] (prob) -- node[above]{\small Conditionnement $\kappa$} (algo);
\draw[arr] (algo) -- node[above]{\small Stabilit\'e} (res);
\draw[arr,dashed,red] (prob) to[bend left=30] node[above]{\small Erreur totale $\leq\kappa\cdot\eps$} (res);
\end{tikzpicture}
\end{center}

\section{Impl\'ementation}

\begin{codeblock}[title=\textbf{Python}]
\begin{minted}{python}
import numpy as np

# Epsilon machine
eps = np.finfo(float).eps
print(f"eps_mach = {eps}")  # 2.220446049250313e-16

# Cancellation catastrophique
x = 1e-8
f_bad = (1 - np.cos(x)) / x**2
f_good = 2 * np.sin(x/2)**2 / x**2
print(f"Mauvais : {f_bad}")   # 0.0
print(f"Bon     : {f_good}")  # 0.49999999...

# Conditionnement d'une matrice
A = np.array([[1, 1], [1, 1.0001]])
print(f"cond(A) = {np.linalg.cond(A):.1f}")  # ~40000
\end{minted}
\end{codeblock}

\begin{codeblock}[title=\textbf{Julia}]
\begin{minted}{julia}
# Epsilon machine
println("eps_mach = ", eps(Float64))

# Cancellation
x = 1e-8
f_bad = (1 - cos(x)) / x^2
f_good = 2sin(x/2)^2 / x^2
println("Mauvais : $f_bad")
println("Bon     : $f_good")

# Conditionnement
using LinearAlgebra
A = [1.0 1.0; 1.0 1.0001]
println("cond(A) = ", cond(A))
\end{minted}
\end{codeblock}

\section{Exercices}

\begin{exercice}[$\star$ -- Epsilon machine]
\'Ecrire un programme qui calcule $\eps_{\mathrm{mach}}$ par dichotomie. Comparer avec \texttt{np.finfo(float).eps}.
\end{exercice}

\begin{exercice}[$\star$ -- Cancellation]
Proposer un calcul stable de $\sqrt{x+1}-\sqrt{x}$ pour $x$ grand. V\'erifier num\'eriquement.
\end{exercice}

\begin{exercice}[$\star\star$ -- Conditionnement]
Calculer le conditionnement du probl\`eme $f(x)=\ln x$ au point $x=1$. Interpr\'eter.
\end{exercice}

\begin{exercice}[$\star\star$ -- Matrice de Hilbert]
Soit $H_n$ la matrice de Hilbert ($h_{ij}=1/(i+j-1)$). Calculer $\cond(H_n)$ pour $n=2,\ldots,15$. Que constate-t-on ?
\end{exercice}

\begin{exercice}[$\star\star\star$ -- Projet]
\'Etudier la stabilit\'e num\'erique de l'\'evaluation du polyn\^ome $p(x)=\sum a_ix^i$ par l'algorithme de Horner vs.\ l'\'evaluation na\"ive. Comparer les erreurs relatives pour des polyn\^omes de degr\'e \'elev\'e.
\end{exercice}

\begin{formules}
\begin{align*}
\fl(x) &= x(1+\delta),\;|\delta|\leq\eps_{\mathrm{mach}} &
\eps_{\mathrm{mach}} &= \tfrac{1}{2}\beta^{1-t} \\
\kappa(f,x) &= |xf'(x)/f(x)| &
\cond(A) &= \|A\|\cdot\|A^{-1}\|
\end{align*}
\end{formules}
